\subsection{Spectrum of an element in a unital algebra}

DEFINITION 1. --- Let A be a unital algebra over K and e its identity element. For every \(x \in \mathrm{A}\) we call the spectrum of \(x\) relative to A the set of \(\lambda \in \mathrm{K}\) such that \(\lambda e - x\) is not invertible.

The spectrum of x will be denoted \(\mathrm{Sp}_{\mathrm{A}}(x)\) , or \(\mathrm{Sp}(x)\) if no confusion results. The complement of \(\mathrm{Sp}_{\mathrm{A}}(x)\) in K is called the resolvent set of x.

Remarks. --- 1) If A = \{0\}, then we have Sp(0) = ∅.

\begin{enumerate}
\def\labelenumi{\arabic{enumi})}
\setcounter{enumi}{1}
\item
  If \(A \neq \{0\}\) , we have \(\mathrm{Sp}(\lambda e) = \{\lambda\}\) for all \(\lambda \in K\) .
\item
  For \(x \in \mathrm{A}\) to be invertible, it is necessary and sufficient that \(0 \notin \operatorname{Sp}(x)\) .
\item
  Let \(R \in K(X)\) be a rational fraction and let \(x \in A\) be an element that is substitutable in R, that is (A, IV, p.~20) there exist P and \(Q \in K[X]\) such that \(R = P/Q\) and \(Q(x)\) is invertible; one can then form the element \(R(x) = P(x) \cdot Q(x)^{-1} = Q(x)^{-1} \cdot P(x)\) of A; it does not depend on the choices of P and Q. We have \(0 \notin Q(\mathrm{Sp}(x))\) , so that every element of \(\mathrm{Sp}(x)\) is substitutable in R.
\end{enumerate}

We have \(\mathrm{R}(\mathrm{Sp}(x)) \subset \mathrm{Sp}(\mathrm{R}(x))\). Indeed, let \(\lambda \in \mathrm{Sp}(x)\) ; there exists a polynomial \(P_{1}\) such that \(\mathrm{R}(\lambda) - \mathrm{R}(\mathrm{X}) = (\lambda - \mathrm{X})\mathrm{P}_{1}(\mathrm{X}) / \mathrm{Q}(\mathrm{X})\) ; then, \(\mathrm{R}(\lambda)e - \mathrm{R}(x) = (\lambda e - x)(\mathrm{P}_{1}(x)/\mathrm{Q}(x))\) , so that \(\mathrm{R}(\lambda) - \mathrm{R}(x)\) is not invertible, hence \(\mathrm{R}(\lambda) \in \mathrm{Sp}(\mathrm{R}(x))\) .

Conversely, suppose that the field K is algebraically closed. Suppose first that R is not constant and let us prove that one has \(\mathrm{Sp}(\mathrm{R}(x)) = \mathrm{R}(\mathrm{Sp}(x))\) . Let \(\mu \in \mathrm{Sp}(\mathrm{R}(x))\) . Since R is not constant, P - \(\mu Q\) is not the zero polynomial; let \(\mu Q - P = \alpha \prod (\lambda_i - X)\) be a decomposition into factors of degree 1, so that \(\mu e - R(x) = \alpha \prod (\lambda_i e - x) Q(x)^{-1}\) . Since \(\mu e - R(x)\) is not invertible, there exists i such that \(\lambda_i e - x\) is not invertible, hence \(\lambda_i \in \mathrm{Sp}(x)\) , then \(\mathrm{R}(\lambda_i) = \mu \in \mathrm{R}(\mathrm{Sp}(x))\) .

When R is constant, the equality \(\mathrm{Sp}(\mathrm{R}(x)) = \mathrm{R}(\mathrm{Sp}(x))\) also holds, provided that \(\mathrm{Sp}(x)\) is nonempty.

\begin{enumerate}
\def\labelenumi{\arabic{enumi})}
\setcounter{enumi}{4}
\item
  Suppose that the algebra A is nonzero. Let \(x \in A\) be a nilpotent element. Let n be an integer such that \(x^{n} = 0\). The spectrum of \(x^{n}\) is reduced to 0, hence so is the spectrum of x by remark 4.
\item
  Let A and B be unital algebras over K and \(\varphi\colon\mathrm{A}\to\mathrm{B}\) a unital morphism. For every \(x\in\mathrm{A}\) , one has \(\mathrm{Sp}_{\mathrm{B}}(\varphi(x))\subset\mathrm{Sp}_{\mathrm{A}}(x)\) .
\item
  Let A be a unital algebra, R its radical (A, VIII, p.~150, def. 2) and let \(\varphi\) be the canonical morphism from A onto \(B = A/R\) . If \(x \in A\) , one has \(\mathrm{Sp}_{\mathrm{B}}(\varphi(x)) = \mathrm{Sp}_{\mathrm{A}}(x)\) . Indeed, it suffices to prove that if \(\varphi(x)\) is invertible in B, then \(x\) is invertible in A. Now, if \(y \in A\) is such that \(\varphi(x)\varphi(y) = \varphi(y)\varphi(x) = \varphi(e)\) , one has \(xy \in e + R\) , \(yx \in e + R\) , hence \(xy\) and \(yx\) are invertible (A, VIII, p.~151, th. 1) and consequently \(x\) is invertible. In particular, if \(x \in R\) , one has \(\mathrm{Sp}_{\mathrm{A}}(x) = \{0\}\) if \(A \neq \{0\}\) .
\item
  Let \((\mathrm{B}_{i})_{i\in\mathrm{I}}\) be a family of unital algebras, with \(\mathrm{B}_{i}=(\mathrm{A}_{i},e_{i})\) for \(i\in I\) . Set \(A=\prod_{i}A_{i}, e=(e_{i})_{i\in\mathrm{I}}\) . Then \((\mathrm{A},e)\) is a unital algebra called the product of the \(B_{i}\) . If \(x=(x_{i})_{i\in\mathrm{I}}\in\mathrm{A}\) , one has \(\mathrm{Sp}_{\mathrm{A}}(x)=\bigcup_{i}\mathrm{Sp}_{\mathrm{A}_{i}}(x_{i})\) .
\end{enumerate}

Examples. --- 1) Let X be a set and let A = K \(^{X}\) be the algebra of functions with values in K defined on X. The spectrum of an element f of A is the set of values of f.

\begin{enumerate}
\def\labelenumi{\arabic{enumi})}
\setcounter{enumi}{1}
\tightlist
\item
  Let A be a unital algebra of finite rank over K. For \(x \in A\) to be invertible, it is necessary and sufficient that the linear map \(y \mapsto xy\) from A to A have nonzero determinant. It follows that the spectrum of x is the set of roots of the characteristic polynomial of x (A, III, p.~110). If A is the algebra of endomorphisms of a vector space V of finite dimension over K, the spectrum of x is therefore the set of eigenvalues of x. This is not always the case when V is of infinite dimension (cf.~I, p.~153, exercise 2).
\end{enumerate}

